\subsection{Error budget of the measurement chain}
\label{sec:budget}

The classification counts already include every error the chain actually made.
A separate budget is still useful, because it shows which contributions are large enough to move a feature and which are not.
Start from the sensor.
The B\&K~4321 charge sensitivity is stated as \(9.8\,\mathrm{pC/g}\) within \(\pm 2\%\).
A two-percent scale error multiplies every acceleration sample on that axis by a factor between \(0.98\) and \(1.02\).
Severity, being an RMS of a quantity linear in acceleration, inherits the same relative error:
\begin{equation}
\frac{\Delta\mathrm{VIB}}{\mathrm{VIB}}\approx 0.02.
\label{eq:senserr}
\end{equation}
On the healthy severity of \(1.051\,\mathrm{mm/s}\) this is about \(0.021\,\mathrm{mm/s}\).
On the eccentric severity of \(3.737\,\mathrm{mm/s}\) it is about \(0.075\,\mathrm{mm/s}\).
The gap between healthy and pedestal severities is \(1.161-1.051=0.110\,\mathrm{mm/s}\).
A two-percent calibration error is a fifth of that gap on the healthy reading and is smaller than the gap.
It does not by itself explain a confusion between those two states, and it is negligible beside the factor \(3.56\) that separates the eccentric state from the healthy state.
After the train-only normalization \eqref{eq:norm}, a pure scale error that is constant for the whole campaign is absorbed into \(u^{(j)}\) and \(v^{(j)}\) of the severity coordinate, provided the same sensor and the same gain are used for every record.
A scale error that changed between the training days and the test days would not be absorbed.
The source does not report a recalibration, and the budget assumes the gain was held fixed.

The acquisition card states a DC accuracy of \(\pm 0.1\%\) of full-scale range.
On the \(\pm 10\,\mathrm{V}\) range this is \(10\,\mathrm{mV}\).
On the \(\pm 100\,\mathrm{mV}\) range it is \(0.1\,\mathrm{mV}\).
Which range was selected is not printed.
The time-domain plots show accelerations of roughly \(7\) to \(9\,\mathrm{m/s}^{2}\), about one \(g\).
At \(9.8\,\mathrm{pC/g}\) that motion produces about \(10\,\mathrm{pC}\) at the peaks.
The voltage that results is the charge divided by the amplifier's feedback capacitance, a setting the source also does not print.
The relevant statement the budget can still make is comparative.
A full-scale error of \(0.1\%\) is twenty times smaller than the sensor's \(2\%\) sensitivity tolerance, so on any range that actually uses a reasonable fraction of its bits the card is not the leading scale error.
Quantization at, for example, a \(16\)-bit converter on a range matched to a few \(g\) sits well below both.
The chain is sensor-limited for amplitude, and it is clock-limited for frequency.

Clock error moves the orders in Table~\ref{tab:bins}.
A sampling-rate error of \(\varepsilon\) moves every apparent frequency by the same relative amount.
The widest order, \(6f_{0}=96\,\mathrm{Hz}\), moves by \(96\varepsilon\).
The half-order acceptance window specified for this exposition is \(4\,\mathrm{Hz}\).
Keeping the true \(6f_{0}\) inside that window requires
\begin{equation}
96\,\abs{\varepsilon}<4
\qquad\Rightarrow\qquad
\abs{\varepsilon}<0.042,
\label{eq:clock}
\end{equation}
a four-percent clock, which any PXI timebase meets by a wide margin.
The binding frequency constraint is not the clock.
It is the mechanical speed.
A one-percent speed drift, already discussed as an envelope of the method, moves \(6f_{0}\) by about \(1\,\mathrm{Hz}\), still inside the window, and moves the supply-to-\(3f_{0}\) gap from \(2.00\,\mathrm{Hz}\) to either \(1.52\,\mathrm{Hz}\) or \(2.48\,\mathrm{Hz}\) according to the sign.
The \(1.22\,\mathrm{Hz}\) bin can still hold them apart.
A speed that was only nominally \(960\,\mathrm{r/min}\), and that wandered by several percent inside a record, would broaden every order and would raise the pedestal and seal-wear confusion first, because those states are identified by which order is largest rather than by a large severity.
The source states a single speed.
The budget treats residual speed wander as an unquantified contribution inside the pedestal errors, not as a number the tables support.

Aliasing is the contribution the budget cannot bound from the printed table.
The sensor responds to \(12\,\mathrm{kHz}\) and the converter samples at \(10\,\mathrm{kHz}\), so energy between \(5\) and \(12\,\mathrm{kHz}\) folds into \(0\) to \(5\,\mathrm{kHz}\) unless an analog filter stops it.
Folded energy lands preferentially in the fine wavelet levels, and it can change which level wins \(k^{\star}\).
If the winning level moves, the coefficient rate changes and the amplitudes \eqref{eq:samples} are read on a different series.
The source plots show \(d_{5}\) winning for the displayed records, which is evidence that, on those records, the mid band carried the energy, and it is not evidence that an anti-alias filter was present.
A repetition of the build should state the filter cutoff next to the \(10\,\mathrm{kHz}\) clock.
Until that cutoff is known, high-frequency folding remains an open term in the budget and a candidate contributor to the two pedestal test errors that every SVM method still makes.

Interference tones that the order sampler does not aim at --- the reported crosstalk at \(28\), \(34\), and \(38\,\mathrm{Hz}\) under the eccentric fault, and the supply line at \(50\,\mathrm{Hz}\) --- are rejected by the sampler's narrow windows and are not rejected by the wavelet energy that chooses the level.
They can change \(k^{\star}\) without entering \(\vect{x}\) as their own coordinates.
That indirect path is the reason a clean mechanical spectrum and a clean scale choice are both part of signal conditioning, which is the improvement the source itself defers to future work.
The budget's conclusion is short.
Calibration and quantization are too small to explain the published confusions.
Speed error at the one-percent level is tolerable for the order windows.
Unfiltered high-frequency energy and in-band interference are large enough, in mechanism, to move the pedestal features, and they are exactly the mechanisms named in the source discussion of that class.

\subsection{An operating procedure for one labeled record}
\label{sec:procedure}

The build is reproducible as a sequence even where some numerical settings remain those of the original run.
One labeled record, from the command to start the pump through the row that enters \(\vect{X}\), is the following procedure.
It is the system of Section~\ref{sec:build} written as steps an operator can tick, not as a new experiment.
\begin{enumerate}
\item Set the loop to the duty in Table~\ref{tab:rig}: \(960\,\mathrm{r/min}\), \(550\,\mathrm{m}^{3}/\mathrm{h}\), \(3\,\mathrm{m}\) head, four-blade impeller, closed circuit.
Confirm \(f_{0}=16\,\mathrm{Hz}\) from a tachometer or from the once-per-revolution line before any fault label is trusted.
\item Install the fault for this record, or install none.
Eccentric axle: displace the balance weight and record the mass and the radius, even though the classifier never sees those two numbers, because a later reader of the severity \(3.737\,\mathrm{mm/s}\) will want them.
Pedestal: seed the bearing defect the campaign has defined and confirm that the pedestal is the vibrating part under the sensor.
Seal wear: use the worn ring the campaign has defined and confirm that the impeller is not rubbing.
Do not combine two of these in one record.
The label set is a partition.
\item Mount the triaxial charge accelerometer on the casing.
Note the direction that will be called radial, because that direction supplies the single severity coordinate.
Check that the charge cables and the amplifier gains are the same ones used for the records already in the training matrix.
\item Set the converter to \(10\,\mathrm{kHz}\), confirm the anti-alias cutoff is below \(5\,\mathrm{kHz}\), and arm a capture of at least \(8192\) samples on all three axes after the pump has been at speed long enough for the flow to be steady.
Store the three series and the tachometer reading with the state label.
\item Compute radial severity by the campaign's stated integrator and by \eqref{eq:vibdisc}.
Run the eight-level db4 bank on each axis, select \(k^{\star}\) by \eqref{eq:kstar} with the finer level winning a tie, and sample the seven orders of \eqref{eq:ordersexact} from the coefficient power spectrum \eqref{eq:cps} through the frequency map \eqref{eq:fmap}.
Refuse any order above the coefficient Nyquist of the winning level.
\item Assemble \eqref{eq:x22}.
If the row is a training row, update nothing about scaling until the full training matrix exists, then fit \eqref{eq:norm}.
If the row is a test row, apply the training ranges already stored.
Do not refit them.
\item When \(180\) training rows are in hand, run the modified leap inside the box \eqref{eq:box} with fitness \eqref{eq:fitness}, retrain the six binary machines on all \(180\) rows at the returned pair, and only then label the \(60\) test rows by vote.
Log the six items of Section~\ref{sec:cost}.
\end{enumerate}
Steps 1--4 are the rig.
Steps 5 and 6 are VWC.
Step 7 is MSFLA-SVM.
A baseline replaces step 7 and does not replace steps 1--6.
That separation is what makes Table~\ref{tab:test} a comparison of classifiers on one feature.

The procedure also says what the operator does not do.
The operator does not choose \(C\) and \(g\) by hand after looking at the test labels.
The operator does not drop the pedestal records because they are harder.
The operator does not lengthen the test set by moving training rows across the split after the search has run.
Each of those moves would raise a number in Table~\ref{tab:test} and would destroy the estimate.
The published counts are informative because, as far as the source text shows, none of those moves was made.

\subsection{The estimate in the unit the experiment controls}
\label{sec:unit}

Table~\ref{tab:errors} restates the published result as error counts, which is the unit in which a one-vector gap is not asked to look like a percentage-point discovery.
Training errors are out of \(180\).
Test errors are out of \(60\).
The per-class test errors follow by subtracting the correct counts of Table~\ref{tab:test} from \(15\).

\begin{table}[H]
\centering
\caption{Error counts. Training errors are out of \(180\) pooled vectors. Test errors are out of \(60\) pooled vectors, and the four class columns are out of \(15\).}
\label{tab:errors}
\begin{tabular}{lrrrrrr}
\toprule
Method & Train & Test & Normal & EAF & BPF & SRWF \\
\midrule
BPNN & 15 & 9 & 2 & 2 & 3 & 2 \\
ACROA-SVM & 5 & 5 & 0 & 2 & 2 & 1 \\
SFLA-SVM & 6 & 5 & 0 & 1 & 2 & 2 \\
MSFLA-SVM & 3 & 4 & 0 & 1 & 2 & 1 \\
\bottomrule
\end{tabular}
\end{table}

Read across the test columns, the modified search matches the better of the other two searches in every state: it matches SFLA-SVM on the eccentric axle (one error), matches both SVM searches on the pedestal (two errors), matches ACROA-SVM on seal wear (one error), and matches every SVM search on the healthy state (no error).
It does not create a new kind of success on the pedestal.
It collects, in one rule, the per-class results that the other searches hold only separately.
That is a fair description of a one-vector pooled gain, and it is the description this estimate keeps.

Read down the pedestal column, no method is in single digits of a different kind: the best pedestal result on the test set is two errors out of fifteen, which is \(13/15\).
The training pedestal column of Table~\ref{tab:train} is consistent with the same ranking.
MSFLA-SVM misses two of forty-five pedestal training vectors, and every other method misses more.
The feature, not the leap, is saturated on this class.
Further optimization of \((C, g)\) can rearrange one error among the easier classes, which is what the sealing-ring column shows, and it has not been observed to buy a pedestal test vector.

The practical statement that closes the estimate is therefore a statement about effort.
On this rig, with this sample, the VWC vector plus a radial SVM already carries the diagnosis from the neural baseline's nine test errors to five.
The modified leap carries it from five to four, by avoiding one seal-wear error that the unmodified leap makes and one eccentric error that the chemical-reaction search makes.
The two pedestal errors that remain are the next experiment: a second sensor, a stated anti-alias filter, or a conditioning step aimed at the site tones, scored on a new split.
They are not a reason to discard the twenty-two coordinates or the leap that reads them.
They are the part of the problem the published counts show is still open.
